%!TEX root = newmain.tex

The definition of $\StipulaToJava{\A}{S}$ follows.

\vspace{-.1cm}

\begin{itemize}
\item[--] $\StipulaToJava{\A}{ \zero} \eqdef \varepsilon$, where
  $\varepsilon$ is the empty string.
%\medskip
\item[--]
  $\StipulaToJava{\A}{E \, \send \, {\x} \; \; S} \eqdef \; \, \x \,
  \mbox{\tt =} \, E \ttsemi \StipulaToJava{\A}{ S}$.
%\medskip
\item[--]
  $\StipulaToJava{\A}{E \, \send \, {\B} \; \; S} \eqdef
  \StipulaToJava{\A}{ S}$. The value of $E$ is irrelevant in our
  examples. If the value of $E$ is relevant, we declare an \emph{ad
    hoc} field, say \stipula{B_transmitted}, for storing the value.
%\medskip
\item[--] $\StipulaToJava{\A}{E \, \lolli \, \h, \h' \; \; S}$. There
  are two cases: (\emph{i}) both $\h$ and $\h'$ are assets of the
  contract or 
  
  (\emph{ii}) $\h$ is a formal parameter and $\h'$ is an
  asset of the contract. The code for case (\emph{i}) and (\emph{ii})
  is ({\tt tmp} is a fresh variable):
  \begin{center}
    {\small
      \begin{tabular}{c@{\qquad\quad}c}
        $
        \begin{array}{l}
          {\tt int \; tmp} \, \mbox{\tt =} \, E \ttsemi 
          \\
          \mbox{\jml{//@ $\;$ assert h-tmp >= 0}} \ttsemi 
          \\
          \h \, \mbox{\tt =} \, \h \mbox{\tt -} {\tt tmp} 
          \ttsemi \h' \, \mbox{\tt =} \, \h' \mbox{\tt +} {\tt tmp} \ttsemi 
          \\
          \StipulaToJava{\A}{ S}
        \end{array}
        $
        &
          $
          \begin{array}{l}
            {\tt int \; tmp} \, \mbox{\tt =} \, E \ttsemi 
            \\
            \mbox{\jml{//@ $\;$ assert h-tmp >= 0  \&\& \A${\unds}$\h'-tmp >= 0}} \ttsemi 
            \\
            \h \, \mbox{\tt =} \, \h \mbox{\tt -} {\tt tmp} 
            \ttsemi \A{\unds}\h' \, \mbox{\tt =} \, \A{\unds}\h' \mbox{\tt -} {\tt tmp} 
            \ttsemi
            \\
            \h' \, \mbox{\tt =} \, \h' \mbox{\tt +} {\tt tmp} \ttsemi 
            \\
            \StipulaToJava{\A}{ S}
          \end{array}
          $
        \\[.2cm]
        Case \emph{(i)} & Case \emph{(ii)}
      \end{tabular}}
  \end{center}

\noindent
% In case {\tt tmp} is a fresh identifier storing the value of the
% expression $E$.
According to the \JML\ \jml{assert} semantics, verification can only
succeed provided that \jml{h-tmp >= 0}, which is exactly what the
{\Stipula} semantics requires. In addition, in case (\emph{ii}), the
formal parameter {\h} represents (part of) the asset {\A}{\tt
  \_}{\h}$'$ that the caller {\A} transfers to the contract upon
invocation. Consequently, the decrease in {\tt h} must be reflected in
the caller's asset rather than in the contract state. This semantic
correspondence motivates the assignment {\tt A{\unds}h$'$ =
  A{\unds}h$'$-tmp;} which correctly accounts for the transfer of
the asset from the caller to the contract.

\medskip

\item[--] $\StipulaToJava{\A}{E \, \lolli \, \h, \B \; \; S}$.  Again,
  there are two cases. Either (\emph{i}) {\h} is an asset or
  (\emph{ii}) it is a formal parameter. The code of cases (\emph{i})
  and (\emph{ii}) is:
  %
  \begin{center}
    {\small
      \begin{tabular}{c@{\qquad\quad}c}
        $\begin{array}{l}
           {\tt int \; tmp} \, \mbox{\tt =} \, E \ttsemi 
           \\
           \mbox{\jml{//@ $\;$ assert h-tmp >= 0}}  \ttsemi 
           \\
           \h \, \mbox{\tt =} \, \h \mbox{\tt -} {\tt tmp} \ttsemi 
           \\
           \B{\unds}\h \, \mbox{\tt =} \, \B{\unds}\h \mbox{\tt +} {\tt tmp} \ttsemi 
           \\
           \StipulaToJava{\A}{ S}
         \end{array}$
        &
          $\begin{array}{l}
             {\tt int \; tmp} \, \mbox{\tt =} \, E \ttsemi 
             \\
             \mbox{\jml{//@ $\;$ assert h-tmp >= 0 \&\& \A-tmp >= 0}}  \ttsemi 
             \\
             \h \, \mbox{\tt =} \, \h \mbox{\tt -} {\tt tmp} 
             \ttsemi \A \, \mbox{\tt =} \, \A \mbox{\tt -} {\tt tmp} \ttsemi 
             \\
             \B \, \mbox{\tt =} \, \B \mbox{\tt +} {\tt tmp} \ttsemi 
             \\
             \StipulaToJava{\A}{ S}
           \end{array}$
        \\[.2cm]
        Case \emph{(i)} & Case \emph{(ii)}
      \end{tabular}}
  \end{center}

\noindent
In case (\emph{ii}), we decrease and increase the assets {\A} and
{\B}, respectively, where {\tt tmp} is a fresh variable. In case
(\emph{i}), we assume that $\h$ is an asset, not a formal
parameter. This allows us to determine what asset of {\tt B} must be
increased by the move (the general case is dealt with \emph{ad hoc}
annotations).

\medskip

\item $\StipulaToJava{\A}{\ifte{E}{S'}{S''} \; S} \eqdef \; \,$
    {\small
      $\begin{array}{l} {\tt if} \, \mbox{\tt (} \, E \, \mbox{\tt )
         \{} \; \StipulaToJava{\A}{S'} \; \mbox{\tt \}}
         \\
         {\tt else} \; \mbox{\tt \{} \; \StipulaToJava{\A}{ S''} \;
         \mbox{\tt \}}
         \\
         \StipulaToJava{\A}{ S}
       \end{array}$
     }
     
\medskip

\noindent
The semantics of the conditional is obvious: we simply use the
corresponding {\Java} statement.
\end{itemize} 

Regarding the shortcuts $\h \lolli \h'$ and $\h \lolli \A$, we
generate optimized {\Java} code, to avoid the declaration of the
auxiliary variable {\tt tmp}:
\[
  \StipulaToJava{\A}{\h \lolli \h'} \eqdef \; {\small \jml{h' = h'+h ; h = 0 ;}}\ \mbox{\textcolor{commentColor}{{\small {\tt // \h \; and \h' assets}}}}
\]
\[
  \StipulaToJava{\A}{\h \lolli \h'} \eqdef \;
  {\small
    \begin{array}{ll} 
      \mbox{\jml{//@ $\;$ assert \A${\unds}$h'-h >= 0}} \ttsemi & 
                                                                   \mbox{\textcolor{commentColor}{{\small {\tt // \h\ formal param., \h' asset}}}}
      \\
      \jml{h' = h'+h ;} & 
      \\
      \jml{A_h' = A_h'-h ;} & 
      \\
      \jml{h = 0 ;} &
    \end{array}
  }
\]
%and 
\[
\begin{array}{l@{\quad}l}
{\rm and} & \StipulaToJava{\A}{\h \lolli \B} \eqdef \; \jml{B_h = B_h+h ; h = 0 ;}\ \mbox{\textcolor{commentColor}{{\small {\tt // \h \; asset}}}} 
\\[.2cm]
& \StipulaToJava{\A}{\h \lolli \B} \eqdef \; 
\begin{array}{l}
\mbox{\jml{//@ $\;$ assert \A-h >= 0}} \ttsemi\ \mbox{\textcolor{commentColor}{{\small {\tt // \h \; formal parameter}}}} 
\\
\jml{B = b+h ; h = 0 ;}
\end{array}
\end{array}
\]

%%% Local Variables:
%%% mode: latex
%%% TeX-master: "newmain"
%%% End:
